% !TeX root = ../Book1.tex
\begin{figure}[htbp]
	\centering
	\begin{tikzpicture}[scale=2,font=\small]
		\draw[GeometryGray] (0,0) circle[radius=1];
		\draw[GeometryGray] (1,0) circle[radius=1];
		\fill[GeometryBlue!8] (0,0)--(1,0)--(0.5,{sqrt(3)/2})--cycle;
		\draw[thick,GeometryBlue] (0,0)--(1,0)--(0.5,{sqrt(3)/2})--cycle;
		\draw[GeometryGray,dashed] (0.5,{-sqrt(3)/2})--(0.5,{sqrt(3)/2});
		\draw[GeometryGray,dashed] (0,0)--(0.5,{-sqrt(3)/2})--(1,0);
		\foreach \x/\y in {0/0,1/0,0.5/0}{\fill (\x,\y) circle[radius=0.025];}
		\fill[GeometryBlue] (0.5,{sqrt(3)/2}) circle[radius=0.025];
		\fill[GeometryBlue] (0.5,{-sqrt(3)/2}) circle[radius=0.025];
		\node[below left] at (0,0) {\(A\)};
		\node[below right] at (1,0) {\(B\)};
		\node[below right] at (0.5,0) {\(M\)};
		\node[above] at (0.5,{sqrt(3)/2}) {\(C_+\)};
		\node[below] at (0.5,{-sqrt(3)/2}) {\(C_-\)};
		\node[right] at (0.5,0.44) {\(h\)};
	\end{tikzpicture}
	\caption[等边三角形作图中的两个交点]{在实欧氏平面中，两圆的共同点位于 \(AB\) 的中垂线上，距中点 \(M\) 的长度为 \(h=\sqrt3\lVert B-A\rVert/2\)。图画出的是实数模型；若改变点集，图中可见的位置未必仍是该模型中的点。}
	\label{g1:fig:equilateral-historical-construction}
\end{figure}
